% ============================================================
%  Chapter 2: Related Work
% ============================================================

\chapter{Related Work}
\label{ch:related}

Our work sits at the intersection of four research traditions: neuroevolution,
modular neural networks, self-organizing and developmental systems, and the
theory of Boolean function learning. This chapter reviews each in turn and
identifies the gap that \modnet{} addresses.

% ------------------------------------------------------------
\section{Neuroevolution: Evolving Neural Network Topologies}
\label{sec:neat}
% ------------------------------------------------------------

The idea that neural network topology can be evolved rather than designed
dates back to the early 1990s. \citet{gruau1994neural} introduced
\emph{Cellular Encoding}, in which a genetic program simultaneously evolves
the architecture, weights, thresholds, and biases of a neural network. The
genotype is a tree of graph-rewriting rules inspired by L-systems; through
development, this tree unfolds into a network of arbitrary size and shape.
Gruau demonstrated that cellular encoding could generate modular neural
networks and used it to evolve a controller for a hexapod robot. This early
work established the principle that complex structure can emerge from a
compact generative description.

The NeuroEvolution of Augmenting Topologies (NEAT) algorithm of
\citet{stanley2002evolving} became the canonical framework for topology-evolving
neuroevolution. NEAT starts from a minimal network and incrementally adds nodes
and connections through structural mutation. It introduces three key ideas:
historical markings to allow meaningful crossover between networks of different
topologies, speciation to protect structural innovation from premature
competition, and complexification from small initial topologies. NEAT was shown
to outperform fixed-topology methods on reinforcement learning benchmarks
\citep{stanley2002efficient} and has since been applied to domains ranging from
autonomous rover navigation \citep{shrestha2025reinforced} to scheduling
problems \citep{neat2024scheduling}.

HyperNEAT \citep{stanley2009hypercube} extended NEAT with an indirect encoding
based on Compositional Pattern-Producing Networks (CPPNs). Instead of directly
encoding each connection, HyperNEAT evolves a CPPN whose output at any pair of
coordinates specifies the weight of the corresponding connection. This
geometric encoding exploits regularity in the substrate and scales to
networks with millions of connections \citep{clune2011performance}. The
Enhanced Hypercube-based Encoding (ES-HyperNEAT) of \citet{risi2012enhanced}
further allowed the placement, density, and connectivity of neurons to be
evolved rather than fixed.

Recent work has focused on accelerating NEAT and extending its reach.
\citet{wang2024tensorized} introduced a tensorization method that transforms
the diverse topologies of NEAT into uniformly shaped tensors, enabling GPU
acceleration and reducing training time by an order of magnitude.
\citet{lyubchev2025revisiting} revisited the original NEAT formalism and
proposed improvements to the mutation and selection operators.

Despite their success, these methods share a common assumption: the network is
organized into layers or a predefined substrate geometry. HyperNEAT's CPPN
operates on a grid; NEAT's complexification adds nodes to hidden layers;
tensorized NEAT pads topologies into fixed-shape tensors. None of these
approaches explores fully layer-free computation in which \emph{any node may
connect to any other node}, including itself, and in which the module structure
itself emerges from the evolutionary process.

% ------------------------------------------------------------
\section{Modular Neural Networks}
\label{sec:modular}
% ------------------------------------------------------------

Modularity is a well-established design principle in software engineering.
\citet{parnas1972modularity} argued that systems should be decomposed into
modules with clean interfaces, hidden implementation details, low coupling
between modules, and high cohesion within modules. This principle has been
remarkably successful in software design \citep{larman2004applying}, and it is
natural to ask whether it applies to neural computation as well.

\citet{alho2024modularity} conducted a systematic review of 86 studies on
modular neural networks and found that nearly two-thirds report improvements
in task accuracy compared to monolithic architectures, with 82\% reporting
benefits in training or inference efficiency. This suggests that modularity
is not merely a design preference but a computationally meaningful structural
property.

In the context of neuroevolution, \citet{khare2005coevolutionary} proposed a
co-evolutionary model in which a population of modules is evolved alongside a
population of complete systems that combine modules. They demonstrated that if
a particular task decomposition is beneficial, it can be discovered through
this mechanism. However, their model requires the designer to specify the
number of modules and the interface between them.

\citet{landassuri2012evolution} investigated the emergence of modularity in
evolved neural networks and found that pure-modular architectures emerge at
low connectivity values, and that the more different two tasks are, the
greater the modularity obtained during evolution. This suggests that modularity
is an emergent property of evolutionary search under certain conditions, not a
design choice imposed from above.

More recent work has examined modularity in the context of reinforcement
learning. \citet{munn2024towards} used NEAT to generate networks of varying
modularity on three RL tasks and found that structural modularity correlated
with performance on tasks that admitted natural task decomposition, but not on
others. This result highlights the importance of measuring modularity in a
task-appropriate way.

Related work in hierarchical modular networks includes \citet{rodriguez2025hierarchical},
who showed that multi-level hierarchical modular networks, when implemented as
recurrent neural network reservoirs, enhance memory capacity, support
multitasking, and give rise to a broader range of temporal dynamics compared
to strictly modular or random networks. The modular growth curriculum of
\citet{2024modulargrowth} similarly showed that gradual modular growth of RNNs
provides advantages for learning increasingly complex tasks.

The key distinction between these works and \modnet{} is that in \modnet{},
modular structure is not merely a measurable property of an evolved network;
it is an explicit design element whose size, count, and internal connectivity
are all under evolutionary control. Modules are not predefined, and
cross-module connections emerge as the task demands.

% ------------------------------------------------------------
\section{Self-Organizing and Developmental Systems}
\label{sec:selforg}
% ------------------------------------------------------------

A parallel line of research investigates neural networks inspired by
biological development. Rather than evolving a fixed structure, these systems
grow, prune, and reorganize themselves through local rules.

\citet{plantec2024evolving} proposed the Lifelong Neural Developmental Program
(LNDP), a class of self-organizing networks capable of synaptic and structural
plasticity in an activity- and reward-dependent manner. Their results
demonstrated that structural plasticity is advantageous in environments
requiring fast adaptation or with non-stationary rewards. The LNDP grows from
a single cell using a Neural Developmental Program, but the growth rules
themselves are evolved, not the topology directly.

The GRN-SO-RC model of \citet{hajizada2024developmental} uses a gene regulatory
network (GRN) to regulate synaptic and structural plasticity in a reservoir
computing framework. The structure of the reservoir self-organizes to adapt to
the specific task through the GRN mechanism. This work is notable for
demonstrating that self-organized reservoir structures can match or exceed the
performance of hand-designed topologies.

The broader relationship between neuroevolution and gene regulatory networks
was reviewed by \citet{2025neuroevolution}, who noted that GRNs are more robust
to mutation and thus more evolvable than direct neuroevolution. This
observation motivates the use of compact, generative encodings rather than
direct encodings of network structure.

\citet{hazan2024self} proposed a biologically inspired developmental algorithm
that grows a functional, layered neural network from a single initial cell.
The growth process is driven by noise and local self-organization, and the
resulting networks achieve competitive performance on standard benchmarks.

Finally, the broader relationship between self-organization and neural
computation was reviewed by \citet{gershenson2010guiding}, who surveyed eight
methods for guiding the self-organization of random Boolean networks toward
the critical dynamical regime. This work is particularly relevant because the
critical regime—poised between order and chaos—is believed to maximize
information processing capacity, and is a natural target for self-organizing
neural systems.

% ------------------------------------------------------------
\section{Biological Inspiration: Criticality and Small-World Topology}
\label{sec:bio}
% ------------------------------------------------------------

Biological neural networks exhibit two structural properties that have been
extensively studied: critical dynamics and small-world topology.

The criticality hypothesis states that the brain operates near the phase
transition between order and disorder, where information transmission is
maximized \citep{beggs2007criticality}. \citet{poil2025network} reviewed the
mechanisms by which this critical state is maintained, identifying
self-organized criticality (SOC) as a leading candidate.
\citet{2026organoid} demonstrated that human forebrain organoids self-organize
to a computationally favorable critical state without external input,
providing empirical support for the self-organization hypothesis. In the
context of modular networks, \citet{2025hierarchicalmodular} showed that
structural modularity tunes mesoscale criticality in biological neuronal
networks.

Small-world topology—characterized by high clustering and short path
lengths—has also been studied extensively. \citet{2024emergence} evolved a
recurrent spiking neural network through neuroevolution and found that
spontaneously evolved topologies exhibit biologically plausible small-world
properties, including densely connected hub nodes, short paths, and long-tailed
degree distributions. \citet{2026evolutionarily} similarly found that
evolutionarily optimized network topologies exhibit sparse connectivity,
small-world organization, and modular architecture, and that these structural
features—not connection density—drive performance advantages over matched
random and Watts–Strogatz baselines.

Scale-free connectivity is a related concept. \citet{mocanu2018scalable}
showed that during training, the hidden neurons of a sparse restricted
Boltzmann machine naturally evolve from an Erd\H{o}s–R\'enyi topology toward a
scale-free one. In the context of Hopfield networks, \citet{2004topology}
showed that scale-free topologies store and retrieve binary patterns more
efficiently than random-diluted networks with the same number of synapses.

These biological and topological findings provide both motivation and
validation for \modnet{}. Our architecture does not enforce small-world or
scale-free structure; we observe that sparsity, cross-module integration, and
recurrence emerge from evolutionary pressure alone, without explicit
regularization for any of these properties.

% ------------------------------------------------------------
\section{Boolean Function Learning and Parity}
\label{sec:boolean}
% ------------------------------------------------------------

The learning of Boolean functions by neural networks has been a foundational
problem since the perceptron. \citet{minsky1969perceptrons} proved that a
single-layer perceptron cannot compute the XOR function, nor the parity
function for any number of inputs. Their proof rests on the fact that
threshold functions are monotone (or anti-monotone) in each variable, and
therefore cannot compute a function as sensitive as parity.

The theoretical hardness of parity for bounded-depth threshold circuits was
subsequently quantified. \citet{impagliazzo1997size} showed that any
threshold circuit of depth $d$ that computes the parity function on $n$
variables must have at least $n^{1 + c\theta^{-d}}$ edges for some constants
$c, \theta > 0$. This is the first superlinear lower bound for an explicit
function that holds for any fixed depth and applies to threshold circuits with
unrestricted weights. \citet{2018averagecase} proved average-case lower
bounds: parity has correlation at most $n^{-\varepsilon_d}$ with depth-$d$
threshold circuits that have at most $n^{1+\varepsilon_d}$ wires. These results
imply that solving parity at scale requires either increasing depth, increasing
width, or introducing recurrence—all of which \modnet{} exploits.

Empirically, several approaches have addressed the parity problem directly.
\citet{iyoda2003parity} introduced a single neuron with a periodic activation
function that can compute $n$-bit parity exactly for any $n$.
\citet{grochowski2007learning} studied the learning of highly non-separable
Boolean functions using constructive feedforward networks, and demonstrated
that the depth-3 architecture of the Boolean inner product is substantially
harder to train than a depth-2 network. More recently, \citet{2024pac}
studied the learnability of Boolean functions on deep neural networks from a
PAC-learning perspective, and found that smaller formulas are harder to learn,
and that underconstrained 3-CNF formulas are more challenging than
overconstrained ones.

\modnet{} addresses the parity problem from a different angle: rather than
designing a specific activation function or training a fixed architecture, we
allow the network to discover a topology that computes parity through
evolution. Our results show that 4-bit parity can be solved exactly with 12
nodes, and 6-bit parity to 93.8\% accuracy with 10 nodes—without any
parity-specific design.

% ------------------------------------------------------------
\section{Sparse and Efficient Neural Computation}
\label{sec:sparse}
% ------------------------------------------------------------

Biological neural networks are highly sparse: the probability that any two
neurons are connected is on the order of $10^{-6}$ \citep{lecun2015deep}.
This sparsity is not accidental; it reflects metabolic and wiring-length
constraints \citep{chklovskii2004wiring, cherniak1994component}.

In artificial neural networks, sparsity has been studied primarily as a
compression technique. The Sparse Evolutionary Training (SET) method of
\citet{mocanu2018scalable} evolves an initial sparse topology into a
scale-free topology during training, allowing the training of networks with
over one million neurons on commodity hardware. \citet{liu2019sparse}
extended this approach to show that vanilla SET-MLP can be implemented from
scratch using just pure Python and SciPy, and that the resulting networks
achieve competitive accuracy on standard benchmarks.
\citet{2019cosine} proposed a cosine-similarity-based selection criterion for
SET that improves the efficiency of the evolutionary topology update.

More recent work has considered pruning from an evolutionary perspective.
\citet{shah2026pruning} modeled parameter groups as populations whose
influence evolves under selection pressure, arguing that sparsity is not
merely a design choice but an emergent property of selection dynamics under
resource constraints. \citet{hershey2025balancing} showed that sparse RNNs can
be balanced with appropriate hyperparameterization, and that this balance
benefits meta-learning.

The relationship between sparsity and modularity is intimate.
\citet{landassuri2012evolution} found that pure-modular architectures emerge
at low connectivity values, suggesting that sparsity and modularity are two
faces of the same underlying structural pressure. In \modnet{}, the pruning
mechanism operates without any explicit sparsity target: nodes and connections
are removed when they become inactive or isolated, and the resulting networks
end up with 27–67\% active-wire density, depending on the task.

% ------------------------------------------------------------
\section{Neural Architecture Search}
\label{sec:nas}
% ------------------------------------------------------------

Neural Architecture Search (NAS) automates the design of network
architectures. Since the influential work of \citet{zoph2017neural},
which formulated NAS as a reinforcement learning problem, the field has grown
rapidly. Modern NAS methods include evolutionary approaches
\citep{real2019regularized}, gradient-based approaches such as DARTS
\citep{liu2019darts}, and one-shot methods \citep{pham2018efficient}.
\citet{ren2024advances} and \citet{ozcelik2026evolutionary} provide recent
surveys.

\modnet{} differs from mainstream NAS in two important respects. First, the
search space is not a discrete set of predefined building blocks (e.g.,
convolutional cells or attention modules); it is the space of all directed
graphs over a fixed number of nodes, with node-level operations drawn from a
small set of weighted-sum-plus-threshold units. Second, the search is
performed by a compact evolutionary algorithm that does not use gradient
information, and the resulting architecture is not a layered stack but a
recurrent graph with modular structure.

% ------------------------------------------------------------
\section{Universal Computation and Neural Cellular Automata}
\label{sec:universal}
% ------------------------------------------------------------

The theoretical foundations of neural computation were established by
\citet{siegelmann1995computational}, who proved that recurrent neural networks
with rational weights are Turing complete. More recently, \citet{perez2021turing}
showed that transformers are Turing complete without external memory, and
\citet{bates2026pytorch} released a library of Turing-complete neural networks
for reproducible experimentation. \citet{stogin2024provably} constructed a
provably stable neural network Turing machine with finite precision and time.

Cellular automata provide a complementary perspective. Rule 110 was proven
Turing complete by \citet{cook2004universality}, and reversible
two-dimensional cellular spaces were shown to support universal computation
by \citet{morita2002universal}. Neural cellular automata, introduced by
\citet{mordvintsev2020growing}, train a cellular automaton to grow and
regenerate a target image from a single seed. \citet{bena2025path} recently
demonstrated that neural cellular automata can achieve universal computation
in a continuous setting.

\modnet{} shares with neural cellular automata the idea that computation
can be local, recurrent, and free of any predefined layer structure. However,
while neural cellular automata operate on a fixed grid with a shared update
rule, \modnet{} allows the topology itself—including the number and
connectivity of nodes—to be evolved.

% ------------------------------------------------------------
\section{Summary and Positioning}
\label{sec:positioning}
% ------------------------------------------------------------

Table~\ref{tab:positioning} summarizes the relationship between \modnet{}
and the most closely related architectures.

\begin{table}[h]
  \centering
  \small
  \caption{Positioning of \modnet{} relative to related architectures.}
  \label{tab:positioning}
  \begin{tabular}{lccccc}
    \toprule
    \textbf{Architecture} & \textbf{Layer-free} & \textbf{Recurrent}
      & \textbf{Modular} & \textbf{Grown} & \textbf{Gradient-free} \\
    \midrule
    NEAT \citep{stanley2002evolving} & No & Optional & No & Yes & Yes \\
    HyperNEAT \citep{stanley2009hypercube} & No & No & No & Yes & Yes \\
    Cellular Encoding \citep{gruau1994neural} & Yes & Yes & Yes & Yes & Yes \\
    Cascade-Correlation \citep{fahlman1990cascade} & No & No & No & Yes & No \\
    LNDP \citep{plantec2024evolving} & Yes & Yes & No & Yes & No \\
    GRN-SO-RC \citep{hajizada2024developmental} & Yes & Yes & No & Yes & No \\
    Neural CA \citep{mordvintsev2020growing} & Yes & Yes & No & Yes & No \\
    \textbf{\modnet{} (this work)} & \textbf{Yes} & \textbf{Yes}
      & \textbf{Yes} & \textbf{Yes} & \textbf{Yes} \\
    \bottomrule
  \end{tabular}
\end{table}

The distinguishing features of \modnet{} are: (i) no predefined layers or
substrate geometry; (ii) recurrence as a first-class architectural primitive;
(iii) modular structure whose size, count, and internal connectivity are all
under evolutionary control; (iv) growth through both addition and pruning;
and (v) no gradient information used at any stage of learning. To the best of
our knowledge, no prior architecture simultaneously satisfies all five of
these constraints.

The remainder of this thesis describes the \modnet{} architecture
(Chapter~\ref{ch:method}), the experimental setup (Chapter~\ref{ch:experiments}),
and the results obtained on a suite of Boolean and continuous tasks
(Chapter~\ref{ch:results}).